\section{Native Backend Experimental Details}
\label{sec:external-native-evidence}

This section documents the uniform scaling cohorts and complete-tree controls reported in the main text. The full artifact also retains prefix-concentrated and clustered sensitivity runs, which are outside these reported cohorts. The same sparse proof task is executed with two external PIR backends. These native-record measurements are separate from the earlier 432-proof SimplePIR and 20-proof VBPIR component checks, and from the Fuel and certificate-mirror TCP studies in Section~\ref{sec:layout-transfer}, where each digest uses eight independent 32-bit SimplePIR chunks. Those earlier measurements remain available under their original configuration; their times, traffic, and bootstrap costs are not pooled with the native 256-bit results. The PBC-routed entry in Table~\ref{tab:comparison_sources} refers to the earlier matching-based adapter; the native PBC adapters used here are specified below.

\subsection{Sources, parameters, and adapters}

The VBPIR implementation follows Mughees and Ren~\cite{vectorbatchpir}, using the \texttt{VBPIR\_PBC} implementation in the TreePIR repository~\cite{treepir}, commit \texttt{930063c5}. The original Client, Server, and PirParams computation is retained; observational getters expose allocated coefficient payloads. The adapter adds original-level proof reconstruction, ciphertext serialization/loading, and public interval routing. PBC retains the official three-candidate placement. Its route table is cleared for each request. The complete-tree controls retain official Java CSA bucket order and validate SubCSA locations, while using a common online directory adapter. Consequently, they compare layouts on one backend without attributing the adapter's client state or routing time to TreePIR's native fast-indexing service.

The VBPIR context uses degree 8192, coefficient-modulus bit lengths $(42,58,58,60)$, a 28-bit plaintext modulus, first dimension 64, and SEAL's \texttt{tc128} context validation. One Client object, one secret key, and its associated evaluation keys serve all buckets through packed queries. This is the main text's shared-state batch realization, whose privacy assumes the underlying encryption's IND-CPA security and VBPIR's evaluation-key circular security; context parameter validation alone is not a protocol-security proof. SimplePIR~\cite{simplepir} uses official commit \texttt{e9020b03}, with unchanged Init, Setup, Query, and Answer routines. Its long-record representation stacks base-$p$ digits vertically. We extend integer input and final reconstruction to 256 bits using exact big integers; modular subtraction, denoising, and rounding follow the original Recover routine. Thus a 32-byte digest requires one query, rather than eight independent scalar databases.

All SimplePIR methods apply \texttt{PickParams} to their public maximum bucket load with record length 256, LWE dimension 1024, and modulus $2^{32}$. The resulting configurations use $p=991$, Gaussian standard deviation 6.4, and 26 digits per record. Buckets within a layout share matrix dimensions; Flat initializes one database and reuses its hint and public matrix for $m$ separate full-database calls. Every call uses fresh private randomness. The SimplePIR PBC wrapper imports the official C++ candidate buckets and record positions and implements the empty-slot-first, random-eviction routing rule in Go. It uses deterministic batch insertion and an independent routing RNG; this does not reproduce the C++ unordered-map eviction trajectory or modify cryptographic randomness. Public $A$ is explicitly expanded and serialized in this experiment; seed compression is a separate available optimization. These settings retain the official parameter policy without asserting identical cryptographic costs across backends.

\subsection{Cohorts, measurements, and authentication}

The uniform cohorts contain 1000, 10000, and 100000 occupied coordinates at original height 128. The three snapshots use seed 2026091103 and nested occupied-coordinate sets. Both backends use byte-identical active digests, values, occupied coordinates, defaults, original levels, and authenticated roots. Each snapshot has a fixed pool of 32 targets; each repeat selects ten without replacement, paired across methods and backends. Targets may recur across repeats. SimplePIR evaluates AB, PBC, Flat-active, and First-fit; VBPIR evaluates AB and PBC. The complete-tree controls have heights 10 and 14 and evaluate AB, PBC, and official CSA under VBPIR. These explicitly reported cohorts comprise 540 uniform measured proofs and 180 complete-tree measured proofs.

Each method/case has three fresh sequential processes, two warmups per process, and ten measured proofs. The reported cohorts therefore contain 72 processes, 720 measured proofs, and 144 warmups. Independent Python verification authenticates all 864 returned proofs, derives required siblings from original heap coordinates, checks actual recovered bytes and bucket positions, and rejects modified siblings, coordinates, and values. A fixed public request envelope includes dummy slots. The statistical unit is a process mean; error bars use sample standard deviations, and paired differences use descriptive 95\% Student-$t$ intervals with two degrees of freedom. These statements concern the stated cohorts, without estimating behavior across arbitrary snapshot distributions.

Execution uses an Intel Core Ultra 5 245K, Ubuntu 24.04, GCC 13.3, and Go 1.22.2 for SimplePIR, with one execution thread. The continuous online timer includes routing, query generation, actual uncompressed serialization/loading, server computation, response serialization/loading, recovery, proof assembly, and root verification. Publisher-oracle checks and negative controls are outside timing. This is local serialized proof processing; network delay is not included.

\subsection{Initialization and stored state}

\begin{table*}[t]
\centering
\fontsize{8.5}{10}\selectfont
\setlength{\tabcolsep}{3pt}
\caption{Uniform-cohort initialization and state. Paired cells list AB (ours)/PBC~\cite{sealpir}; native placement follows VBPIR~\cite{vectorbatchpir}. Each method has three complete processes. $I$ is mean setup time in seconds. $H$ is the database-dependent hint, $A/K$ is expanded SimplePIR $A$ or VBPIR evaluation-key payload, $D$ is the serialized public directory, and $S$ is total serialized initialization components, including relevant framing and DBinfo. $H,A/K,D,S,F$ are MiB; combined peak RSS is GiB. $F$ is the common-metadata cache reference. VBPIR has no SimplePIR-style downloaded hint. Byte quantities are deterministic; RSS and setup values are means.}
\label{tab:external-initialization}
% Generated from frozen JSON by build_external_supplement_table.py.
% Each paired numerical cell is AB / PBC.
\begin{tabular}{lr r r r r r r r}
\toprule
Backend & Leaves & $I$ (s) & $H$ & $A/K$ & $D$ & $S$ & $F$ & RSS (GiB)\\
\midrule
VBPIR~\cite{vectorbatchpir} & 1,000 & 0.33/0.37 & -- & 37.57/37.57 & 0.11/0.17 & 37.68/37.74 & 0.14 & 0.45/0.45 \\
VBPIR~\cite{vectorbatchpir} & 10,000 & 0.69/1.02 & -- & 37.57/37.57 & 1.07/1.68 & 38.64/39.25 & 1.38 & 0.45/0.45 \\
VBPIR~\cite{vectorbatchpir} & 100,000 & 4.38/8.29 & -- & 37.57/37.57 & 10.69/16.79 & 48.26/54.36 & 13.74 & 1.16/1.64 \\
\midrule
SimplePIR~\cite{simplepir} & 1,000 & 0.08/0.19 & 3.96/8.12 & 2.64/6.72 & 0.15/0.18 & 6.75/15.03 & 0.20 & 1.21/1.21 \\
SimplePIR~\cite{simplepir} & 10,000 & 0.36/0.78 & 12.09/26.41 & 11.22/24.17 & 1.46/1.86 & 24.77/52.44 & 1.97 & 1.21/1.21 \\
SimplePIR~\cite{simplepir} & 100,000 & 2.04/4.71 & 42.66/87.75 & 39.21/88.00 & 15.30/19.68 & 97.17/195.43 & 20.17 & 1.20/1.21 \\
\bottomrule
\end{tabular}

\end{table*}

Table~\ref{tab:external-initialization} separates initialization time, public state, and memory. Setup includes input parsing and backend preparation. SimplePIR loads precomputed bucket maps for every layout, excluding external tree and layout construction. VBPIR likewise excludes external AB/CSA construction but includes its runtime PBC bucket construction. SimplePIR computes actual database-dependent hints. Its received serialized state contains hints, expanded $A$, directory metadata, and DBinfo. VBPIR initialization also includes evaluation keys uploaded to the server; this bidirectional accounting is not persistent client memory. Combined peak RSS includes both roles, retained inputs, and validation buffers.

The cache reference contains all active digests plus the same backend-specific common metadata encoding, excluding bucket maps and cryptographic keys. The digest-only sizes are 0.061, 0.610, and 6.103\,MiB, respectively. Neither reference is a minimum-client-state claim. For a fixed snapshot, initialization and online costs must both be considered when comparing private retrieval with caching. Per-process matrices, state components, input hashes, schedules, and the source-selection manifest are retained for reproduction.
